\documentclass[10pt,twocolumn,a4paper]{article}
\usepackage[top=2cm, bottom=2cm, left=1.5cm, right=1.5cm]{geometry}
\usepackage{fontspec}
\usepackage[portuguese, provide=*]{babel}
\usepackage{enumitem}
\setlist[itemize]{label=-}

\usepackage{amsmath}
\usepackage{amssymb}
\usepackage{booktabs}
\usepackage{float}
\usepackage{hyperref}
\usepackage{listings}
\usepackage{graphicx}
\usepackage{xcolor}

% Configuração para blocos de código
\definecolor{codegreen}{rgb}{0,0.6,0}
\definecolor{codegray}{rgb}{0.5,0.5,0.5}
\definecolor{codepurple}{rgb}{0.58,0,0.82}
\definecolor{backcolour}{rgb}{0.95,0.95,0.92}

\lstdefinestyle{mystyle}{
    backgroundcolor=\color{backcolour},
    commentstyle=\color{codegreen},
    keywordstyle=\color{magenta},
    numberstyle=\tiny\color{codegray},
    stringstyle=\color{codepurple},
    basicstyle=\ttfamily\scriptsize,
    breakatwhitespace=false,
    breaklines=true,
    captionpos=b,
    keepspaces=true,
    numbers=left,
    numbersep=5pt,
    showspaces=false,
    showstringspaces=false,
    showtabs=false,
    tabsize=2,
    language=Python
}
\lstset{style=mystyle}

% --- Marcadores de pendência ---
\newcommand{\pend}[1]{\textbf{\textcolor{red}{[PREENCHER: #1]}}}
\newcommand{\pp}{\textbf{\textcolor{red}{[PREENCHER]}}}
\newcommand{\figpend}[2]{%
  \fbox{\parbox[c][#1][c]{0.94\columnwidth}{\centering\textcolor{red}{\textbf{[FIGURA PENDENTE]}}\\ \footnotesize #2}}}

\title{\vspace{-1.5cm}\textbf{Busca Local e Algoritmo Genético para o Sequenciamento de Manutenção de Ativos: uma Análise Comparativa}}
\author{
    \textbf{Rafael Campello Soares, Arthur Rafael Silva Teixeira}\\
    \textbf{Roberto Júnior Cassimiro Campolina, Lucas Araújo Campos Szuster}\\
    \textit{Universidade Federal de Minas Gerais (UFMG) -- Departamento de Engenharia Elétrica}\\
    \textit{Disciplina: Fundamentos de Inteligência Artificial -- Trabalho Prático 2}\\
    \textit{Belo Horizonte, MG, Brasil}
}
\date{}

\begin{document}

\maketitle

% Inclusion das seções separadas
\noindent\textbf{\textit{Resumo}}\; Este trabalho compara três métodos de otimização --- \textit{greedy descent} com \textit{random restart}, \textit{simulated annealing} e algoritmo genético --- aplicados ao sequenciamento de manutenção preventiva de ativos de uma concessionária de energia, formulado como a minimização do atraso ponderado total sobre permutações. Os métodos são avaliados em três instâncias ($n=12$, $n=50$ e $n=150$) sob um orçamento idêntico de avaliações da função objetivo e com execuções independentes.
\pend{resumir os principais resultados quantitativos e a conclusão central}

\vspace{0.3cm}
\noindent\textbf{\textit{Palavras-chave}}\; busca local; \textit{simulated annealing}; algoritmo genético; otimização combinatória; sequenciamento; atraso ponderado total.

\vspace{0.5cm}
%======================================================================
\section{Introdução}
%======================================================================

Uma concessionária de distribuição de energia precisa executar a manutenção preventiva de $n$ ativos (religadores, transformadores de distribuição e chaves seccionadoras) com uma única equipe, que parte de uma base e atende um ativo por vez. Cada ativo possui posição, duração de serviço, prazo regulatório e um peso de criticidade que representa o custo por hora de atraso. A ordem de visita determina tanto o tempo gasto em deslocamentos quanto os atrasos incorridos, e o objetivo é encontrar a ordem que minimiza o atraso ponderado total.

O espaço de busca possui $n!$ soluções; para $n=150$, esse número tem mais de 260 algarismos, o que inviabiliza qualquer enumeração. Trata-se de um problema de otimização combinatória sem restrições rígidas: toda permutação é viável e há apenas uma função de custo a minimizar. Dois critérios gulosos antagônicos --- visitar o ativo mais próximo ou o mais urgente --- não produzem, isoladamente, boas soluções, o que torna o problema adequado para comparar métodos de busca mais elaborados.

O objetivo do trabalho é comparar, em um problema prático de engenharia, heurísticas de busca local e uma metaheurística populacional, observando na prática o equilíbrio entre \textit{exploração} (visitar regiões diferentes do espaço de busca) e \textit{intensificação} (refinar soluções em uma região promissora). Os três algoritmos são:

\begin{itemize}
  \item \textbf{(A)} \textit{greedy descent} com \textit{random restart};
  \item \textbf{(B)} \textit{simulated annealing};
  \item \textbf{(C)} algoritmo genético (\textit{genetic algorithm}, GA).
\end{itemize}

A comparação segue um protocolo único: o recurso contabilizado é o número de avaliações da função $h$, e todos os algoritmos recebem o mesmo orçamento por instância (Seção~\ref{sec:protocolo}). A pergunta central, a ser respondida com dados na Seção~\ref{sec:resultados}, é qual a relação entre cada algoritmo e a estrutura do problema --- em particular, se a recombinação do GA consegue combinar bons trechos de sequência de dois pais.

\textbf{Organização.} A Seção~\ref{sec:modelagem} apresenta a modelagem formal. A Seção~\ref{sec:algoritmos} descreve a arquitetura do código e os três algoritmos. A Seção~\ref{sec:protocolo} detalha a metodologia experimental. A Seção~\ref{sec:resultados} apresenta os resultados e a análise comparativa, e a Seção~\ref{sec:conclusoes} traz as conclusões.
%======================================================================
\section{Modelagem do Problema}\label{sec:modelagem}
%======================================================================

\subsection{Dados da instância}

Seja $\mathcal{A}=\{1,\dots,n\}$ o conjunto de ativos e $0$ o índice da base. Para cada ativo $j\in\mathcal{A}$ são dados a posição $(x_j,y_j)$, a duração do serviço $p_j$ (horas), o prazo $d_j$ (hora-limite de conclusão) e o peso de criticidade $w_j$ (custo por hora de atraso). A base tem posição $(x_0,y_0)$ e $p_0=0$. O tempo de deslocamento entre dois pontos $i$ e $j$ é
\begin{equation}
  t_{ij}=\operatorname{round}\!\left(\sqrt{(x_i-x_j)^2+(y_i-y_j)^2}\right),
  \label{eq:tij}
\end{equation}
com arredondamento para o inteiro mais próximo. Esse arredondamento é necessário para que os valores de $h$ sejam comparáveis entre grupos. A equipe não retorna à base ao final.

\subsection{Espaço de estados, ações e função objetivo}

\begin{itemize}
  \item \textbf{Estado (solução candidata):} uma permutação $\pi=(\pi_1,\dots,\pi_n)$ de $\mathcal{A}$, que representa a ordem de visita. O espaço de busca é $\mathcal{S}=S_n$, com $|\mathcal{S}|=n!$.
  \item \textbf{Ações:} os movimentos definidos pelos operadores de vizinhança (Seção~\ref{sec:vizinhancas}). Cada movimento transforma uma permutação em outra.
  \item \textbf{Restrições:} não há restrições rígidas, pois toda permutação é viável. Os prazos $d_j$ são restrições brandas (\textit{soft constraints}), penalizadas na função objetivo. O problema é de otimização, e não de satisfação.
  \item \textbf{Instante de conclusão:} com $\pi_0=0$ e $C_{\pi_0}=0$,
  \begin{equation}
    C_{\pi_k}=C_{\pi_{k-1}}+t_{\pi_{k-1}\pi_k}+p_{\pi_k}.
    \label{eq:conclusao}
  \end{equation}
  \item \textbf{Função objetivo (\textit{fitness}), a minimizar:} o atraso ponderado total
  \begin{equation}
    h(\pi)=\sum_{j=1}^{n} w_j\cdot\max(0,\,C_j-d_j).
    \label{eq:h}
  \end{equation}
  \item \textbf{Problema:} $\pi^\star=\arg\min_{\pi\in S_n}h(\pi)$.
\end{itemize}

Na terminologia do algoritmo genético, o \textbf{cromossomo} é a permutação completa, o \textbf{gene} é uma posição da sequência, o \textbf{alelo} é o ativo que a ocupa e a \textbf{aptidão} é $h(\pi)$. Como $h$ mede custo, o GA busca o cromossomo de menor aptidão.

\subsection{Propriedades úteis}

O custo acumulado $h$ é não decrescente ao longo do prefixo da permutação: acrescentar o $(k{+}1)$-ésimo ativo nunca reduz o atraso já acumulado pelos $k$ primeiros. Além disso, o instante $C_{\pi_k}$ depende apenas do prefixo $(\pi_1,\dots,\pi_k)$, o que permite avaliar um movimento recalculando somente o sufixo afetado (Seção~\ref{sec:incremental}).

\subsection{Instâncias}

Os arquivos \texttt{manutencao\_*.txt} contêm $n$, as coordenadas da base e, para cada ativo, $(x_j,y_j,p_j,d_j,w_j)$. A Tabela~\ref{tab:instancias} resume as instâncias e os valores de referência fornecidos no enunciado. O valor da instância pequena é o ótimo global, obtido por busca exata. Os demais são os melhores valores conhecidos, obtidos pelos professores com orçamento muito maior, e servem como alvo, não como meta obrigatória.

\begin{table}[H]
\centering
\caption{Instâncias e valores de referência de $h$.}
\label{tab:instancias}
\small
\begin{tabular}{lrrrl}
\toprule
Instância & $n$ & $\sum_j p_j$ & Ref. $h$ & Natureza\\
\midrule
pequena & 12 & 55 & 1\,875 & ótimo global\\
intermediária & 50 & 239 & 16\,332 & melhor conhecido\\
grande & 150 & 684 & 60\,417 & melhor conhecido\\
\bottomrule
\end{tabular}
\end{table}

\subsection{Operadores de vizinhança}\label{sec:vizinhancas}

Os três operadores clássicos geram sucessores de uma permutação (Figura do enunciado, exemplo com $\pi=(3,7,1,5,2,8)$):

\begin{itemize}
  \item \textbf{Swap (troca):} troca os ativos nas posições $i<j$. Há $\binom{n}{2}=\frac{n(n-1)}{2}$ sucessores.
  \item \textbf{Insertion (realocação):} remove o ativo da posição $i$ e o reinsere na posição $j\neq i$. Há $n(n-1)$ movimentos, mas a realocação entre posições adjacentes coincide com um \textit{swap} adjacente e é contada duas vezes; o número de sucessores distintos é $(n-1)^2$.
  \item \textbf{2-opt (inversão):} inverte o segmento entre as posições $i<j$. Há $\frac{n(n-1)}{2}$ sucessores, e a inversão de um segmento de tamanho 2 coincide com um \textit{swap} adjacente.
\end{itemize}

A Tabela~\ref{tab:vizinhanca} apresenta o tamanho de cada vizinhança nas três instâncias. Os operadores não são equivalentes: cada um induz um conjunto distinto de sucessores e, portanto, uma superfície adaptativa diferente para o mesmo problema.

\begin{table}[H]
\centering
\caption{Tamanho da vizinhança induzida em função de $n$.}
\label{tab:vizinhanca}
\small
\begin{tabular}{lcrrr}
\toprule
Operador & Fórmula & $n=12$ & $n=50$ & $n=150$\\
\midrule
Swap & $n(n-1)/2$ & 66 & 1\,225 & 11\,175\\
Insertion & $(n-1)^2$ & 121 & 2\,401 & 22\,201\\
2-opt & $n(n-1)/2$ & 66 & 1\,225 & 11\,175\\
\bottomrule
\end{tabular}
\end{table}

Na instância grande, uma varredura completa da vizinhança de \textit{swap} consome 11\,175 avaliações, cerca de 3{,}7\% do orçamento de 300\,000. Esse custo condiciona o comportamento das variantes de seleção do sucessor (Seção~\ref{sec:greedy}).
%======================================================================
\section{Algoritmos}\label{sec:algoritmos}
%======================================================================

\subsection{Arquitetura}

Os três algoritmos são genéricos e interagem com o problema apenas por uma interface comum (Tabela~\ref{tab:interface}). Essa separação permite executar as 9 combinações (3 algoritmos $\times$ 3 instâncias) sem alterar o código dos algoritmos e centralizar a contagem de avaliações em um único ponto: cada chamada a $h$ incrementa um contador único, inclusive a reavaliação da solução corrente.

\begin{table}[H]
\centering
\caption{Interface entre algoritmos e problema.}
\label{tab:interface}
\small
\begin{tabular}{p{3.2cm}p{4.6cm}}
\toprule
Método & Papel\\
\midrule
\texttt{solucao\_aleatoria()} & gera uma permutação inicial\\
\texttt{h(s)} & fitness (minimização)\\
\texttt{sucessores(s)} & itera sobre todos os sucessores\\
\texttt{sucessor\_aleatorio(s)} & sorteia um sucessor\\
\texttt{cruzamento(s1,s2)} & recombinação\\
\texttt{mutacao(s)} & mutação\\
\bottomrule
\end{tabular}
\end{table}

\pend{linguagem/versão, bibliotecas, gerador aleatório e política de sementes (uma semente por execução, registrada em disco)}

\subsection{Avaliação incremental}\label{sec:incremental}

Avaliar $h$ do zero custa $O(n)$. Ao aplicar um movimento cuja menor posição afetada é $i$, o prefixo $(\pi_1,\dots,\pi_{i-1})$ permanece inalterado, de modo que $C_{\pi_{i-1}}$ e o custo acumulado até $i-1$ podem ser reaproveitados a partir de vetores de prefixo mantidos para a solução corrente. Basta recalcular o sufixo, a um custo $O(n-i+1)$. Isso vale para \textit{swap}, \textit{insertion} e 2-opt: mesmo que os travessias sejam simétricas, os instantes $C_j$ dentro do trecho alterado mudam e precisam ser recalculados. A avaliação incremental não altera a contagem de avaliações, apenas o tempo de parede.

\pend{vizinhança(s) com avaliação incremental implementada, descrição do teste automático de concordância com a avaliação completa e seu resultado}

O ganho de tempo medido é reportado na Seção~\ref{sec:res-incremental}.

\subsection{(A) \textit{Greedy descent} com \textit{random restart}}\label{sec:greedy}

A partir de uma permutação aleatória, o algoritmo move-se repetidamente para um sucessor que melhora $h$, até atingir um \textbf{ótimo local}: uma permutação sem sucessor melhor na vizinhança adotada. Ao atingir o ótimo local, aplica-se \textit{random restart}: recomeça-se de nova permutação aleatória, preservando a melhor solução já vista, até esgotar o orçamento $B$. As variantes de seleção do sucessor são:

\begin{itemize}
  \item \textbf{\textit{Best improvement}:} varre toda a vizinhança e escolhe o melhor sucessor;
  \item \textbf{\textit{First improvement}:} aceita o primeiro sucessor que melhora $h$;
  \item \textbf{Sucessor aleatório com aceitação de melhora:} sorteia um único sucessor e o aceita somente se melhorar $h$.
\end{itemize}

\pend{variantes efetivamente implementadas (mínimo duas) e a variante adotada nos experimentos principais, com justificativa}

Sob orçamento fixo, as variantes diferem fortemente em passos de melhoria por avaliação: o \textit{best improvement} paga uma vizinhança inteira por passo, enquanto o \textit{first improvement} e a variante aleatória pagam em média menos. A Seção~\ref{sec:res-t2} reporta quantos passos de melhoria cada variante deu dentro do orçamento. Para a variante aleatória, o critério de ótimo local não é determinístico; é preciso definir o critério de reinício.

\pend{critério de detecção de ótimo local/reinício da variante aleatória, se implementada}

\subsection{(B) \textit{Simulated annealing}}\label{sec:sa}

Seja $\pi$ a solução corrente e $\pi'$ um sucessor sorteado ao acaso. Se $h(\pi')\le h(\pi)$, aceita-se $\pi'$; caso contrário, aceita-se com probabilidade
\begin{equation}
  P=\exp\!\left(-\frac{h(\pi')-h(\pi)}{T}\right),
  \label{eq:gibbs}
\end{equation}
a distribuição de Gibbs. Cada sucessor examinado consome uma avaliação.

\textbf{\textit{Annealing schedule}.} Adota-se resfriamento geométrico $T\leftarrow\alpha T$. Como cada iteração consome uma avaliação e o orçamento é $B$, impor $T_f=T_0\,\alpha^{\,B/m}$ com atualização a cada $m$ iterações fornece
\begin{equation}
  \alpha=\left(\frac{T_f}{T_0}\right)^{m/B}.
  \label{eq:alpha}
\end{equation}
Assim, a temperatura chega a $T_f$ exatamente ao esgotar o orçamento, e o resfriamento é ajustado a $B$ de cada instância.

\textbf{Calibração de $T_0$ e $T_f$.} A escala de $T$ depende da magnitude típica de $h(\pi')-h(\pi)$, que varia de dezenas a milhares e muda com a instância; portanto, os parâmetros são calibrados \textit{por instância}. O procedimento adotado deve ser:

\begin{enumerate}
  \item amostrar um conjunto de pares $(\pi,\pi')$ aleatórios por instância e estimar a distribuição de $\Delta=h(\pi')-h(\pi)$ para $\Delta>0$;
  \item escolher $T_0$ tal que a taxa de aceitação inicial de movimentos piores seja próxima de um alvo alto (quase \textit{random walk});
  \item escolher $T_f$ tal que a taxa de aceitação final seja praticamente nula (quase \textit{greedy descent});
  \item verificar a escolha pelo gráfico da taxa de aceitação ao longo da execução.
\end{enumerate}

\pend{alvos de aceitação, tamanho da amostra de calibração e valores finais de $T_0$, $T_f$, $\alpha$ e $m$ por instância (Tabela~\ref{tab:par-sa})}

\begin{table}[H]
\centering
\caption{Parâmetros do \textit{simulated annealing} por instância.}
\label{tab:par-sa}
\small
\begin{tabular}{lcccc}
\toprule
Instância & $T_0$ & $T_f$ & $\alpha$ & Vizinhança\\
\midrule
pequena       & \pp & \pp & \pp & \pp\\
intermediária & \pp & \pp & \pp & \pp\\
grande        & \pp & \pp & \pp & \pp\\
\bottomrule
\end{tabular}
\end{table}

\textbf{Taxa de aceitação.} Registra-se, em janelas de avaliações, a fração de movimentos \textit{piores} aceitos. Uma têmpera bem calibrada começa quase como \textit{random walk} e termina quase como \textit{greedy descent}; um gráfico que permanece alto indica $T$ alta demais, e um que cai a zero logo no início indica $T$ baixa demais.

\subsection{(C) Algoritmo genético}\label{sec:ga}

O GA mantém uma população de $N$ cromossomos e repete: seleção, cruzamento com probabilidade $P_c$, mutação com probabilidade $P_m$ e elitismo, até esgotar o orçamento de avaliações. A população inicial é uma amostra aleatória de $N$ permutações.

\textbf{Seleção e conversão custo$\to$aptidão.} Como $h$ é de minimização, o operador de seleção deve favorecer valores \textit{menores}. No torneio de tamanho $k$, basta escolher o indivíduo de menor $h$ entre $k$ sorteados, sem conversão. Na roleta, é necessária uma aptidão $f(\pi)\ge 0$ crescente com a qualidade, por exemplo $f(\pi)=h_{\max}-h(\pi)$ (com $h_{\max}$ o pior valor da geração) ou $f(\pi)=1/(1+h(\pi))$.

\pend{operador de seleção adotado, tamanho do torneio ou fórmula de conversão custo$\to$aptidão efetivamente usada, e justificativa}

\textbf{Cruzamento.} Cruzamentos clássicos de um ou múltiplos pontos duplicariam ativos e gerariam cromossomos inválidos; é necessário um operador que preserve a permutação (OX, PMX ou cruzamento por arestas).

\pend{operador de cruzamento adotado (OX, PMX ou arestas) e justificativa}

A título de ilustração, o \textit{Order Crossover} (OX) opera como segue, com $n=8$, pais $P_1=(1\,2\,3\,4\,5\,6\,7\,8)$ e $P_2=(3\,7\,5\,1\,6\,8\,2\,4)$ e pontos de corte antes da posição 4 e depois da posição 6:

\begin{enumerate}
  \item copia-se o segmento de $P_1$ entre os cortes para o filho: $(\_\,\_\,\_\,4\,5\,6\,\_\,\_)$;
  \item percorre-se $P_2$ a partir da posição seguinte ao segundo corte (posição 7), de forma circular: $2,4,3,7,5,1,6,8$;
  \item descartam-se os ativos já presentes no segmento ($4,5,6$), restando $2,3,7,1,8$;
  \item preenchem-se as posições livres a partir da posição 7, também circularmente: posições $7,8,1,2,3$ recebem $2,3,7,1,8$;
  \item o filho é $(7\,1\,8\,4\,5\,6\,2\,3)$, que é uma permutação válida.
\end{enumerate}

\textbf{Mutação.} Aplica-se uma pequena taxa de mutação para inserir diversidade e favorecer a exploração. O ponto de partida é a \textit{swap mutation}; os demais operadores da vizinhança também são admissíveis.

\pend{operador de mutação adotado}

\textbf{Elitismo e parâmetros.} O elitismo preserva os $e$ melhores indivíduos de cada geração. A Tabela~\ref{tab:par-ga} reúne os parâmetros. Cada indivíduo novo avaliado consome uma avaliação, e o algoritmo para exatamente em $B$.

\begin{table}[H]
\centering
\caption{Parâmetros do algoritmo genético.}
\label{tab:par-ga}
\small
\begin{tabular}{lccc}
\toprule
Parâmetro & pequena & interm. & grande\\
\midrule
$N$ (população) & \pp & \pp & \pp\\
$P_c$ & \pp & \pp & \pp\\
$P_m$ & \pp & \pp & \pp\\
Elitismo ($e$) & \pp & \pp & \pp\\
Tamanho do torneio & \pp & \pp & \pp\\
\bottomrule
\end{tabular}
\end{table}

\pend{como os parâmetros foram escolhidos e testados}

\textbf{Diversidade.} A cada geração registra-se o número de cromossomos distintos e a distância média entre pares de indivíduos, definida pela distância de Hamming entre posições $d(\pi,\sigma)=|\{k:\pi_k\neq\sigma_k\}|$. A queda dessas medidas diagnostica a convergência prematura.
%======================================================================
\section{Protocolo Experimental}\label{sec:protocolo}
%======================================================================

\subsection{Orçamento e execuções}

Todos os algoritmos gastam o mesmo recurso: o número de avaliações de $h$. Cada sucessor examinado pela busca local conta 1, e cada indivíduo avaliado pelo GA conta 1. Os algoritmos param exatamente ao atingir o orçamento (Tabela~\ref{tab:orcamento}).

\begin{table}[H]
\centering
\caption{Orçamento e execuções independentes.}
\label{tab:orcamento}
\small
\begin{tabular}{lrr}
\toprule
Instância & Orçamento $B$ & Execuções\\
\midrule
pequena ($n=12$) & 50\,000 & 30\\
intermediária ($n=50$) & 150\,000 & 30\\
grande ($n=150$) & 300\,000 & 10\\
\bottomrule
\end{tabular}
\end{table}

\subsection{Métricas e estatística}

Para cada combinação algoritmo $\times$ instância, registra-se o melhor $h$ de cada execução e calculam-se média, desvio-padrão, mínimo e mediana. O gap relativo ao ótimo na instância pequena é $\text{gap}=(h-h^\star)/h^\star$, com $h^\star=1875$. Registram-se também o tempo de parede por execução e o número de avaliações.

As comparações par a par usam o teste não paramétrico de Mann--Whitney \cite{mann1947}, bilateral, com nível de significância $\alpha_s=0{,}05$. Como são três comparações por instância, sugere-se correção de Holm.

\pend{teste e correção para múltiplas comparações efetivamente usados e a biblioteca empregada}

Médias próximas com dispersões grandes não sustentam conclusão; as afirmações de superioridade no texto vêm acompanhadas do $p$-valor.

\subsection{Curvas de convergência}

Para cada instância, registra-se o melhor $h$ encontrado em função do número de avaliações gastas, em uma grade fixa de pontos, e plota-se a mediana das execuções por algoritmo. Para o GA, acrescenta-se a curva da média da população; a distância entre as duas curvas indica a perda de diversidade.

\subsection{Ambiente de execução}

\pend{descrição do hardware (CPU, memória), sistema operacional e versões do interpretador/compilador usados nas medições de tempo}

\subsection{Reprodutibilidade}

O código acompanha um \texttt{README} com os comandos exatos para reproduzir cada tabela e figura. Os resultados brutos de todas as execuções são gravados em arquivos \texttt{.csv}.

\pend{endereço do repositório ou nome do arquivo .zip e os comandos do README para cada tabela e figura (ver Anexo A)}
%======================================================================
\section{Resultados e Análise}\label{sec:resultados}
%======================================================================

\subsection{T1: validação da função de avaliação}

A função $h$ foi validada com os dois valores do enunciado para \texttt{manutencao\_pequena.txt}: a permutação identidade $\pi=(1,\dots,12)$ deve resultar em $h(\pi)=4366$, e o ótimo global conhecido é $h^\star=1875$. Uma heurística gulosa razoável (menor acréscimo imediato de atraso) alcança 3168, cerca de 69\% acima do ótimo.

\begin{table}[H]
\centering
\caption{Saída do programa nos casos de verificação.}
\label{tab:t1}
\small
\begin{tabular}{lrr}
\toprule
Caso & Esperado & Obtido\\
\midrule
Identidade, instância pequena & 4\,366 & \pp\\
Ótimo global, instância pequena & 1\,875 & \pp\\
\bottomrule
\end{tabular}
\end{table}

\pend{saída literal do programa para os dois casos; opcionalmente o exemplo com $n=4$ (custos 39 e 22)}

\subsection{T2: comparação das vizinhanças}\label{sec:res-t2}

Usando \textit{greedy descent} com \textit{random restart} e o mesmo orçamento, comparam-se \textit{swap}, \textit{insertion} e 2-opt nas três instâncias.

\begin{table*}[!t]
\centering
\caption{Comparação das vizinhanças no \textit{greedy descent} com \textit{random restart} (média $\pm$ desvio-padrão do melhor $h$ por execução).}
\label{tab:t2}
\small
\begin{tabular}{llcccc}
\toprule
Instância & Vizinhança & Melhor $h$ (média $\pm$ dp) & Passos de melhoria & Reinícios & Tempo por movimento (s)\\
\midrule
pequena & Swap      & \pp & \pp & \pp & \pp\\
pequena & Insertion & \pp & \pp & \pp & \pp\\
pequena & 2-opt     & \pp & \pp & \pp & \pp\\
\midrule
intermediária & Swap      & \pp & \pp & \pp & \pp\\
intermediária & Insertion & \pp & \pp & \pp & \pp\\
intermediária & 2-opt     & \pp & \pp & \pp & \pp\\
\midrule
grande & Swap      & \pp & \pp & \pp & \pp\\
grande & Insertion & \pp & \pp & \pp & \pp\\
grande & 2-opt     & \pp & \pp & \pp & \pp\\
\bottomrule
\end{tabular}
\end{table*}

\pend{resultados do experimento T2}

\textbf{Discussão e escolha da vizinhança.} \pend{qual vizinhança conduz a melhores ótimos locais, qual é mais barata por movimento com avaliação incremental, discussão dos passos de melhoria de cada variante de seleção sob o orçamento (em especial na instância grande, onde uma varredura de \textit{swap} custa 11\,175 avaliações) e justificativa da vizinhança adotada nos demais experimentos}

\subsection{Ganho da avaliação incremental}\label{sec:res-incremental}

\begin{table}[H]
\centering
\caption{Tempo de avaliação: completa vs.\ incremental.}
\label{tab:incremental}
\small
\begin{tabular}{lccc}
\toprule
Instância & Completa (s) & Incremental (s) & Ganho\\
\midrule
pequena       & \pp & \pp & \pp\\
intermediária & \pp & \pp & \pp\\
grande        & \pp & \pp & \pp\\
\bottomrule
\end{tabular}
\end{table}

\pend{medições de tempo e comentário sobre o ganho observado por tamanho de instância}

\subsection{T3: resultados principais}

\begin{table*}[!t]
\centering
\caption{Resultados dos três algoritmos (melhor $h$ por execução). O gap é calculated apenas na instância pequena, com $h^\star=1875$.}
\label{tab:principal}
\small
\begin{tabular}{llcccccc}
\toprule
Instância & Algoritmo & Média $\pm$ dp & Mediana & Mínimo & Gap & Tempo médio (s) & Avaliações\\
\midrule
pequena & Greedy + restart & \pp & \pp & \pp & \pp & \pp & 50\,000\\
pequena & Sim.\ annealing  & \pp & \pp & \pp & \pp & \pp & 50\,000\\
pequena & Alg.\ genético   & \pp & \pp & \pp & \pp & \pp & 50\,000\\
\midrule
intermediária & Greedy + restart & \pp & \pp & \pp & --- & \pp & 150\,000\\
intermediária & Sim.\ annealing  & \pp & \pp & \pp & --- & \pp & 150\,000\\
intermediária & Alg.\ genético   & \pp & \pp & \pp & --- & \pp & 150\,000\\
\midrule
grande & Greedy + restart & \pp & \pp & \pp & --- & \pp & 300\,000\\
grande & Sim.\ annealing  & \pp & \pp & \pp & --- & \pp & 300\,000\\
grande & Alg.\ genético   & \pp & \pp & \pp & --- & \pp & 300\,000\\
\bottomrule
\end{tabular}
\end{table*}

\pend{valores da Tabela~\ref{tab:principal} a partir dos resultados em disco}

\textbf{Boxplots.} A Figura~\ref{fig:box} compara os três algoritmos em cada instância.

\begin{figure*}[!t]
\centering
\figpend{3.5cm}{boxplots do melhor $h$ por algoritmo, um painel por instância (arquivo fig\_boxplots.pdf)}
\caption{Boxplots do melhor $h$ encontrado por execução, por instância e algoritmo.}
\label{fig:box}
\end{figure*}

\textbf{Curvas de convergência.} A Figura~\ref{fig:conv} apresenta a mediana do melhor $h$ em função das avaliações, e, para o GA, a curva da média da população.

\begin{figure*}[!t]
\centering
\figpend{3.8cm}{curvas de convergência por instância (arquivo fig\_convergencia.pdf)}
\caption{Curvas de convergência (mediana das execuções). A distância entre a curva do melhor e a da média da população do GA indica a perda de diversidade.}
\label{fig:conv}
\end{figure*}

\textbf{Testes estatísticos.}

\begin{table}[H]
\centering
\caption{$p$-valores do teste de Mann--Whitney (bilateral).}
\label{tab:pvalores}
\small
\begin{tabular}{lccc}
\toprule
Comparação & pequena & interm. & grande\\
\midrule
Greedy vs.\ SA & \pp & \pp & \pp\\
Greedy vs.\ GA & \pp & \pp & \pp\\
SA vs.\ GA     & \pp & \pp & \pp\\
\bottomrule
\end{tabular}
\end{table}

\pend{p-valores (e correção, se aplicável) e leitura de quais diferenças são estatisticamente significativas}

\textbf{Tempo e avaliações.} A Tabela~\ref{tab:principal} apresenta o tempo médio por execução ao lado do número de avaliações. O orçamento é igual entre algoritmos, mas o custo por avaliação não é: operadores como o cruzamento, a avaliação incremental e a manutenção de estruturas da população fazem o tempo de parede diferir mesmo com avaliações idênticas.

\pend{discussão da diferença entre as métricas (avaliações vs.\ tempo de parede) com base nos tempos medidos}

\subsection{Discussão}

\textbf{Ótimos locais.}
\pend{número de ótimos locais distintos do \textit{greedy descent} na instância pequena; em quantas execuções cada algoritmo atingiu $h^\star=1875$; relação com a discussão de ótimo local $\times$ ótimo global}

\textbf{Exploração $\times$ intensificação.} Em cada algoritmo, identificam-se os seguintes mecanismos candidatos; a análise deve apoiá-los em dados (Seção~\ref{sec:protocolo}):

\begin{itemize}
  \item \textit{Greedy + restart}: a descida intensifica; o \textit{restart} explora.
  \item \textit{Simulated annealing}: a aceitação de piores (alta $T$) explora; o resfriamento leva à intensificação.
  \item GA: o cruzamento e a mutação exploram; a seleção e o elitismo intensificam.
\end{itemize}

\pend{evidências nos dados de onde o equilíbrio foi bem ou mal ajustado: taxa de aceitação do SA (gráfico), diversidade do GA, número de reinícios do \textit{greedy}}

\begin{figure}[H]
\centering
\figpend{3.2cm}{taxa de aceitação do SA ao longo da execução (arquivo fig\_aceitacao.pdf)}
\caption{Taxa de aceitação de movimentos piores no \textit{simulated annealing}.}
\label{fig:aceitacao}
\end{figure}

\begin{figure}[H]
\centering
\figpend{3.2cm}{diversidade da população do GA por geração (arquivo fig\_diversidade.pdf)}
\caption{Evolução da diversidade da população do algoritmo genético.}
\label{fig:diversidade}
\end{figure}

\textbf{Efeito da escala.}
\pend{o ranking muda entre a instância pequena e a grande? Hipótese (se mudar) ou evidência estatística de estabilidade (se não mudar)}

\textbf{Quando o algoritmo genético compensa?}
\pend{existe instância em que o GA supera significativamente as duas abordagens de busca local? Citar os $p$-valores}

\subsubsection{Experimento adicional: a recombinação combina bons trechos?}

\textbf{Hipótese.} O GA só traz vantagem sobre a busca local se a recombinação conseguir juntar bons trechos de sequência de dois pais em um filho melhor, isto é, se a codificação por permutação preservar blocos úteis (por exemplo, subsequências de ativos com prazos próximos visitadas em boa ordem). Se o cruzamento produzir filhos em geral piores que ambos os pais, o GA funciona como uma busca aleatória guiada pela mutação e pela seleção.

\textbf{Experimento proposto.} Duas medidas complementares sobre o mesmo orçamento e as mesmas instâncias:

\begin{enumerate}
  \item \textbf{Ablação:} comparar o GA completo com uma versão sem cruzamento ($P_c=0$, apenas mutação, seleção e elitismo), com o mesmo $N$ e $B$, usando Mann--Whitney sobre o melhor $h$ final.
  \item \textbf{Qualidade dos filhos:} em cada geração, registrar a fração de filhos com $h$ menor que o de ambos os pais, menor que o do melhor pai e maior que o de ambos, e como essas frações evoluem à medida que a diversidade cai.
\end{enumerate}

\begin{table}[H]
\centering
\caption{Experimento adicional: efeito do cruzamento.}
\label{tab:extra}
\small
\begin{tabular}{lccc}
\toprule
Métrica & pequena & interm. & grande\\
\midrule
Melhor $h$, GA completo (média) & \pp & \pp & \pp\\
Melhor $h$, GA sem cruzamento (média) & \pp & \pp & \pp\\
$p$-valor (Mann--Whitney) & \pp & \pp & \pp\\
Filhos melhores que ambos os pais (\%) & \pp & \pp & \pp\\
Filhos piores que ambos os pais (\%) & \pp & \pp & \pp\\
\bottomrule
\end{tabular}
\end{table}

\pend{resultados do experimento e a interpretação: a hipótese foi sustentada ou refutada pelos dados?}
%======================================================================
\section{Conclusões}\label{sec:conclusoes}
%======================================================================

\pend{síntese dos resultados, lições sobre exploração e intensificação, conclusão sobre a recombinação, limitações e trabalhos futuros}

% Bibliografia
\begin{thebibliography}{9}

\bibitem{poole2023}
D.~Poole and A.~Mackworth, \textit{Artificial Intelligence: Foundations of Computational Agents}, 3rd~ed. Cambridge University Press, 2023. Cap.~4.

\bibitem{eiben2003}
A.~E. Eiben and J.~E. Smith, \textit{Introduction to Evolutionary Computing}. Springer-Verlag, 2003.

\bibitem{gaspar2013}
A.~Gaspar-Cunha, R.~H.~C. Takahashi, and C.~H. Antunes (Orgs.), \textit{Manual de Computação Evolutiva e Metaheurística}. Editora UFMG, 2013.

\bibitem{russell2016}
S.~Russell and P.~Norvig, \textit{Artificial Intelligence: A Modern Approach}, 3rd~ed. Pearson, 2016. Cap.~4.

\bibitem{mann1947}
H.~B. Mann and D.~R. Whitney, ``On a test of whether one of two random variables is stochastically larger than the other,'' \textit{Annals of Mathematical Statistics}, vol.~18, no.~1, pp.~50--60, 1947.

\end{thebibliography}

%======================================================================
\appendix
\section{Reprodução dos resultados}
%======================================================================

\begin{table}[H]
\centering
\caption{Mapa de reprodução (tabelas e figuras $\to$ comandos).}
\label{tab:repro}
\small
\begin{tabular}{p{2.6cm}p{5.2cm}}
\toprule
Item & Comando\\
\midrule
Tabela~\ref{tab:t1} (T1) & \pp\\
Tabela~\ref{tab:t2} (T2) & \pp\\
Tabela~\ref{tab:incremental} & \pp\\
Tabela~\ref{tab:principal} & \texttt{python main.py <instância> <algoritmo>}, uma vez por combinação\\
Tabela~\ref{tab:pvalores} & \pp\\
Tabela~\ref{tab:extra} & \pp\\
Figura~\ref{fig:box} & \pp\\
Figura~\ref{fig:conv} & \pp\\
Figura~\ref{fig:aceitacao} & \pp\\
Figura~\ref{fig:diversidade} & \pp\\
\bottomrule
\end{tabular}
\end{table}

\end{document}